% R15 component; only input paths and enumerated location/grammar prose are edited.
% Source: revisions/2026-10-04-r15/manuscript/fees.tex; commit 1cb3cc9efe135966ec228dfaf51c6841a6dfc97c.
% Generated by code/integrate.py; edit extraction rules there.
% Reviewed source: revisions/2026-10-04-r14/manuscript/theory.tex, line 76, commit f5021cefa71492babfcbfa580e0c984f59a9de26
\subsection{When is state-dependent management worth its fee?}\label{sec:r12fees}
To give the policy differences a decision interpretation, consider a proportional resource cost $\tau\in[0,1)$ for administering a state-dependent consumption rule. Gross withdrawals from capital remain $m_i e^{Y_i}$, and the service absorbs $\tau m_i e^{Y_i}$. Household consumption is $(1-\tau)m_i e^{Y_i}$. Production, gross depletion, the gross-withdrawal adjustment cost, and terminal capital are unchanged. This is a self-financed resource experiment rather than an externally financed utility supplement. It changes the specified fee parameter, not the model used to train or benchmark the gross policy.

Let $A_T=\int_0^T e^{-\rho t}dt$. Log felicity gives the exact identity
\begin{equation}\label{eq:r12fee}
 J_{\mathcal P}^{\tau}(a)-J_{\mathcal P}(m_0)
 =J_{\mathcal P}(a)-J_{\mathcal P}(m_0)+A_T\log(1-\tau).
\end{equation}
Hence a positive certified improvement $L_a$ supports every fee $\tau$ satisfying
\[
 \tau<1-\exp(-L_a/A_T).
\]
 A negative upper endpoint after the fee rules out an improvement for that candidate under the confidence event; an interval spanning zero is inconclusive. The fee grid is fixed before the final simulation. We do not equate wall-clock cost to a monetary fee or call these deliberately stylized balance sheets a calibrated household population. The experiment instead states precisely which adoption decisions are supported by the verified utility differences.
